%==============================================================
%  lecons/algebre/jacobi_exp.tex
%  Formule det(exp A) = exp(tr A)
%==============================================================

\lecontitle{Formule \texorpdfstring{$\det(e^A) = e^{\mathrm{tr}(A)}$}{det(exp A) = exp(tr A)}}{Algèbre linéaire — Déterminant et trace de l'exponentielle}

%=============================================================
%  § 1 — ÉNONCÉ
%=============================================================
\secnum{navyblue}{1}{Énoncé}

\begin{tcolorbox}[thmbox]
  \textbf{Théorème (formule de Jacobi–Liouville).}%
  \index{formule de Jacobi}\index{det(exp A) = exp(tr A)}
  Pour tout $A \in \mathcal{M}_n(\mathbb{K})$ ($\mathbb{K} = \mathbb{R}$ ou $\mathbb{C}$)~:
  \[
    \boxed{\det\!\bigl(e^A\bigr) = e^{\mathrm{tr}(A)}}.
  \]
\end{tcolorbox}

\begin{significationintuitive}
Le déterminant mesure la dilatation des volumes~; la trace en est la version
\emph{infinitésimale}, puisque $\mathrm{tr}(B)$ est la dérivée de
$M \mapsto \det(M)$ au point $M = I$ dans la direction $B$ (c'est précisément le
développement au premier ordre ci-dessous). La formule exprime alors que
l'exponentielle de matrices transporte la structure \emph{additive} de la trace
en structure \emph{multiplicative} du déterminant, exactement comme l'identité
scalaire $e^{x+y} = e^x e^y$ transforme l'addition des réels en multiplication.
\end{significationintuitive}

\rubriq{Deux ingrédients clés}

\begin{mdframed}[style=exbar]
  La preuve repose sur deux propriétés élémentaires~:
  \begin{enumerate}
    \item \textbf{Multiplicativité du déterminant}~:
          $\det(MN) = \det(M)\det(N)$.
    \item \textbf{Développement au premier ordre}~:
          pour $\varepsilon \to 0$,
          \[
            \det(I + \varepsilon B) = 1 + \varepsilon\,\mathrm{tr}(B) + O(\varepsilon^2).
          \]
  \end{enumerate}
\end{mdframed}

\rubriq{Preuve du développement au premier ordre}

\begin{mdframed}[style=proofbar]
  Par définition, $\det(I + \varepsilon B) = \sum_{\sigma \in \mathfrak{S}_n}
  \operatorname{sgn}(\sigma)\prod_{i=1}^n (I+\varepsilon B)_{i,\sigma(i)}$.

  Pour $\sigma = \mathrm{id}$~: $\prod_i(1 + \varepsilon B_{ii})
  = 1 + \varepsilon\,\mathrm{tr}(B) + O(\varepsilon^2)$.

  Pour $\sigma \neq \mathrm{id}$~: $\sigma$ possède au moins deux points non fixes
  $(i,j)$ avec $i \neq j$, donc le produit comporte au moins deux facteurs
  $\varepsilon B_{i\sigma(i)}$ et contribue en $O(\varepsilon^2)$.

  Ainsi $\det(I + \varepsilon B) = 1 + \varepsilon\,\mathrm{tr}(B) + O(\varepsilon^2)$.
  \hfill$\blacksquare$
\end{mdframed}

\decsep{}

%=============================================================
%  § 2 — PREUVE PRINCIPALE : ÉQUATION DIFFÉRENTIELLE
%=============================================================
\secnum{navyblue}{2}{Preuve par équation différentielle}

\begin{mdframed}[style=proofbar]
  Posons $f : \mathbb{R} \to \mathbb{K}^*$, $f(t) = \det(e^{tA})$.

  \medskip
  \textbf{Étape 1 — $f$ est un homomorphisme de $(\mathbb{R},+)$ dans
  $(\mathbb{K}^*,\times)$.}

  Pour tous $t, s \in \mathbb{R}$, les matrices $tA$ et $sA$ commutent,
  donc $e^{(t+s)A} = e^{tA}e^{sA}$. La multiplicativité du déterminant donne
  \[
    f(t+s) = \det(e^{(t+s)A}) = \det(e^{tA})\,\det(e^{sA}) = f(t)\,f(s).
  \]

  \medskip
  \textbf{Étape 2 — Calcul de $f'(0)$.}

  \[
    f'(0) = \lim_{h \to 0}\frac{f(h) - f(0)}{h}
           = \lim_{h \to 0}\frac{\det(e^{hA}) - 1}{h}.
  \]
  Or $e^{hA} = I + hA + O(h^2)$, donc par le développement au premier ordre~:
  \[
    \det(e^{hA}) = \det(I + hA + O(h^2)) = 1 + h\,\mathrm{tr}(A) + O(h^2).
  \]
  Ainsi $f'(0) = \mathrm{tr}(A)$.

  \medskip
  \textbf{Étape 3 — $f$ vérifie une équation différentielle linéaire.}

  Fixons $t \in \mathbb{R}$. Grâce à l'équation fonctionnelle de l'étape~1,
  pour tout $h$~:
  \[
    f(t+h) = f(t)\,f(h).
  \]
  En soustrayant $f(t)$ et en divisant par $h$, puis en faisant $h \to 0$, on
  obtient (l'étape~2 garantissant que $f$ est dérivable en $0$)~:
  \[
    f'(t) = \lim_{h\to 0}\frac{f(t+h)-f(t)}{h}
          = f(t)\,\lim_{h\to 0}\frac{f(h)-1}{h}
          = f'(0)\,f(t) = \mathrm{tr}(A)\,f(t).
  \]
  Ainsi $f$ est solution de l'équation différentielle linéaire
  $f'(t) = \mathrm{tr}(A)\,f(t)$ avec $f(0)=1$. Cet argument est valable
  indifféremment pour $\mathbb{K}=\mathbb{R}$ et $\mathbb{K}=\mathbb{C}$, $f$
  étant à valeurs dans $\mathbb{K}^*$. La solution est unique~:
  $f(t) = e^{t\,\mathrm{tr}(A)}$.

  \medskip
  \textbf{Conclusion.}\enspace En $t = 1$~:
  $\det(e^A) = f(1) = e^{f'(0)} = e^{\mathrm{tr}(A)}$. \hfill$\blacksquare$
\end{mdframed}

\decsep{}

%=============================================================
%  § 3 — PREUVE ALTERNATIVE : FORME DE JORDAN
%=============================================================
\secnum{navyblue}{3}{Preuve alternative via la forme de Jordan}

\begin{mdframed}[style=proofbar]
  On travaille sur $\mathbb{C}$~; le cas réel s'en déduit par densité ou par
  restriction.

  \medskip
  Soit $A = P\,J\,P^{-1}$ la décomposition de Jordan,
  $J = \mathrm{diag}(J_{m_1}(\lambda_1), \ldots, J_{m_r}(\lambda_r))$.
  Par conjugaison~: $e^A = P\,e^J\,P^{-1}$, donc
  \[
    \det(e^A) = \det(e^J)
    = \prod_{i=1}^r \det\!\bigl(e^{J_{m_i}(\lambda_i)}\bigr).
  \]

  Pour un bloc de Jordan $J_m(\lambda) = \lambda I_m + N$~:
  \[
    e^{J_m(\lambda)} = e^\lambda\,e^N = e^\lambda(I_m + N + \tfrac{N^2}{2!}+\cdots),
  \]
  qui est triangulaire supérieure avec $e^\lambda$ sur toute la diagonale.
  Donc $\det(e^{J_m(\lambda)}) = {(e^\lambda)}^m = e^{m\lambda}$.

  En multipliant sur tous les blocs~:
  \[
    \det(e^A) = \prod_{i=1}^r e^{m_i \lambda_i}
              = \exp\!\Bigl(\sum_{i=1}^r m_i \lambda_i\Bigr)
              = e^{\mathrm{tr}(J)}
              = e^{\mathrm{tr}(A)},
  \]
  car $\mathrm{tr}(PJP^{-1}) = \mathrm{tr}(J) = \sum_i m_i\lambda_i$. \hfill$\blacksquare$
\end{mdframed}

\decsep{}

%=============================================================
%  § 4 — COROLLAIRES
%=============================================================
\secnum{navyblue}{4}{Corollaires}

\begin{tcolorbox}[thmbox]
  \begin{enumerate}
    \item \textbf{Positivité réelle.}\enspace%
          \index{det(exp A) > 0}
          Pour tout $A \in \mathcal{M}_n(\mathbb{R})$~:
          $\det(e^A) = e^{\mathrm{tr}(A)} > 0$.
          L'exponentielle réelle a son image dans
          $\mathrm{GL}_n^+(\mathbb{R}) = \{M \mid \det M > 0\}$.

    \item \textbf{Déterminant unité.}\enspace{}
          $\det(e^A) = 1 \iff \mathrm{tr}(A) = 0$.
          En particulier~:
          \[
            \exp : \mathfrak{sl}_n(\mathbb{K}) \to \mathrm{SL}_n(\mathbb{K}),
          \]
          où $\mathfrak{sl}_n = \{A \mid \mathrm{tr}(A)=0\}$ est l'algèbre de
          Lie\index{algèbre de Lie} de $\mathrm{SL}_n$.

    \item \textbf{Inversibilité de $e^A$.}\enspace{}
          ${(e^A)}^{-1} = e^{-A}$, que l'on retrouve~:
          $\det(e^A)\,\det(e^{-A}) = e^{\mathrm{tr}(A)}\,e^{-\mathrm{tr}(A)} = 1$.

    \item \textbf{Formule de Liouville.}\enspace%
          \index{formule de Liouville}
          Soit $X' = A(t)\,X$ un système différentiel ($A$ pouvant dépendre de $t$)
          et $W(t) = \det\bigl(x_1(t),\ldots,x_n(t)\bigr)$ le wronskien
          de $n$ solutions. Alors~:
          \[
            W(t) = W(0)\,\exp\!\!\int_0^t \mathrm{tr}(A(s))\,ds.
          \]
  \end{enumerate}
\end{tcolorbox}

\rubriq{Preuve du corollaire 4 — formule de Liouville}

\begin{mdframed}[style=proofbar]
  Posons $\Phi(t)$ la matrice dont les colonnes sont $x_1(t),\ldots,x_n(t)$
  (matrice fondamentale). Elle vérifie $\Phi' = A(t)\Phi$.

  En dérivant $W(t) = \det\Phi(t)$ et en utilisant la multilinéarité du
  déterminant par rapport aux \emph{lignes}~: en notant $L_1,\ldots,L_n$
  les lignes de $\Phi(t)$,
  \[
    W'(t) = \sum_{i=1}^n
      \det\begin{pmatrix}L_1\\ \vdots\\ L_i'\\ \vdots\\ L_n\end{pmatrix},
  \]
  où, dans le $i$-ème terme, seule la $i$-ème ligne est dérivée.
  Or la relation $\Phi' = A(t)\Phi$ se lit ligne par ligne~: en notant
  $A(t)=(a_{k\ell}(t))$, on a $L_i' = \sum_{k=1}^n a_{ik}(t)\,L_k$
  (les coefficients sont ici bien ceux de $A$), d'où
  \[
    \det\begin{pmatrix}L_1\\ \vdots\\ L_i'\\ \vdots\\ L_n\end{pmatrix}
      = \sum_{k=1}^n a_{ik}(t)\,
        \det\begin{pmatrix}L_1\\ \vdots\\ L_k \;\text{(place $i$)}\\ \vdots\\ L_n\end{pmatrix}.
  \]
  Pour $k \neq i$, le déterminant possède deux lignes égales
  (la ligne $L_k$ figure déjà en position $k$ et est répétée en position $i$)~:
  il est donc nul. Seul le terme $k=i$ survit, et il vaut
  $a_{ii}(t)\,\det\Phi(t) = a_{ii}(t)\,W(t)$. En sommant sur $i$~:
  \[
    W'(t) = \sum_{i=1}^n a_{ii}(t)\,W(t) = \mathrm{tr}(A(t))\,W(t).
  \]
  Cette équation différentielle linéaire scalaire s'intègre en
  $W(t) = W(0)\exp\!\bigl(\int_0^t\mathrm{tr}(A(s))\,ds\bigr)$.
  $\blacksquare$
\end{mdframed}

\decsep{}

%=============================================================
%  § 5 — PIÈGES ET EXERCICES
%=============================================================
\secnum{exgreen}{5}{Pièges et exercices}

\rubriq{Pièges}

\begin{mdframed}[style=cexbar]
  \textbf{La réciproque est fausse~: $\det M > 0$ n'implique pas $M \in \mathrm{Im}(\exp)$
  sur $\mathbb{R}$.}\enspace{}
  Comme montré dans la leçon sur la surjectivité,
  $\bigl(\begin{smallmatrix}-1&1\\0&-1\end{smallmatrix}\bigr)$ a déterminant $1$
  mais n'est pas dans l'image de $\exp:\mathcal{M}_2(\mathbb{R})\to\mathrm{GL}_2(\mathbb{R})$.

  \smallskip\noindent
  \textbf{Ne pas confondre $\det(e^A)$ et $e^{\det(A)}$.}\enspace{}
  Ces deux quantités sont généralement différentes~: pour
  $A = \begin{pmatrix}1&0\\0&2\end{pmatrix}$,
  $\det(e^A) = e^3$ mais $e^{\det(A)} = e^2$.

  \smallskip\noindent
  \textbf{La formule de Liouville suppose $A$ \emph{quelconque}, pas seulement constante.}\enspace{}
  Lorsque $A$ est constante, $\Phi(t) = e^{tA}$ et
  $W(t) = \det(e^{tA}) = e^{t\,\mathrm{tr}(A)}$, ce qui redonne la formule principale.
\end{mdframed}

\vspace{6pt}
\exolabel{}

\begin{mdframed}[style=exobar]
  \begin{enumerate}
    \item Montrer que $\det(e^A) = e^{\mathrm{tr}(A)}$ en utilisant
          uniquement la densité des matrices diagonalisables dans
          $\mathcal{M}_n(\mathbb{C})$ et la continuité des deux membres.

    \item Soit $A \in \mathcal{M}_n(\mathbb{R})$ antisymétrique ($A^T = -A$).
          Montrer que $\det(e^A) = 1$ et que $e^A$ est une matrice
          orthogonale.

    \item Pour $A = \begin{pmatrix}a&b\\c&d\end{pmatrix}$, vérifier
          directement que $\det(e^A) = e^{a+d}$ en calculant $e^A$
          dans le cas $A$ diagonalisable, puis dans le cas $A$ nilpotente,
          enfin dans le cas général (décomposition de Dunford).

    \item Soit $X' = AX$ avec $A \in \mathcal{M}_2(\mathbb{R})$ constante et
          $W(0)=1$. Exprimer $W(t)$ et montrer que le signe de $W(t)$
          est constant.

    \item (\textit{Défi}) Établir la formule de Jacobi différentielle~:
          pour $M : \mathbb{R} \to \mathcal{M}_n(\mathbb{R})$ inversible et
          différentiable,
          \[
            \frac{d}{dt}\det M(t) = \det(M(t))\,\mathrm{tr}\!\bigl({M(t)}^{-1}M'(t)\bigr).
          \]
          En déduire une troisième preuve de $\det(e^A) = e^{\mathrm{tr}(A)}$.
  \end{enumerate}
\end{mdframed}

\leconfooter{%
  C.\ G.\ J.\ Jacobi (1841) —
  J.\ Liouville (formule du wronskien, 1838)%
}
